\subsubsection{Color Register (COLR)}

Address offset: \textbf{0x28}
\newline
Reset value: 0x0000 0000

\begin{figure}[H]
    \centering
    \begin{bytefield}[bitwidth=\bflenhalfword, endianness=big]{16}
        \bitheader[lsb=16]{16-31} \\
        \bitbox{12}{Reserved} &
        \bitbox{4}{CB} \\
        \bitbox[]{12}{} &
        \bitbox[]{4}{R/W}
    \end{bytefield}
\end{figure}

\begin{figure}[H]
    \centering
    \begin{bytefield}[bitwidth=\bflenhalfword, endianness=big]{16}
        \bitheader{0-15} \\
        \bitbox{4}{Res} &
        \bitbox{4}{CG} &
        \bitbox{4}{Res} &
        \bitbox{4}{CR} \\
        \bitbox[]{4}{} &
        \bitbox[]{4}{R/W} &
        \bitbox[]{4}{} &
        \bitbox[]{4}{R/W}
    \end{bytefield}
\end{figure}

\begin{itemize}
    \item \textbf{Bits 31:20 - Reserved}
    \item \textbf{Bits 19:16 - CB: Color Blue}\\
    Der blaue Farbanteil des in VDATR gesetzten Zeichens im in ADDRR ausgewählten Pixelblock.
    Wird übernommen, wenn WD in CTRL auf 1 gesetzt wird.
    \item \textbf{Bits 15:12 - Reserved}
    \item \textbf{Bits 11:8 - CG: Color Green}\\
    Der grüne Farbanteil des in VDATR gesetzten Zeichens im in ADDRR ausgewählten Pixelblock.
    Wird übernommen, wenn WD in CTRL auf 1 gesetzt wird.
    \item \textbf{Bits 7:4 - Reserved}
    \item \textbf{Bits 3:0 - CR: Color Red}\\
    Der rote Farbanteil des in VDATR gesetzten Zeichens im in ADDRR ausgewählten Pixelblock.
    Wird übernommen, wenn WD in CTRL auf 1 gesetzt wird.
\end{itemize}
